% Part: first-order-logic
% Chapter: tableaux
% Section: rules-and-proofs

\documentclass[../../../include/open-logic-section]{subfiles}

\begin{document}

\iftag{FOL}
      {\olfileid{fol}{tab}{rul}}
      {\olfileid{pl}{tab}{rul}}

\olsection{规则与分析表}

分析表是对一个语句在一个结构中可能为真或为假的各种方式的系统考察。
分析表的基本构件是带号公式：语句加上一个真值“标记”，即 $\True$ 或~$\False$。
这些带号公式排列成一棵（向下生长的）树。

\begin{defn}
  一个\emph{带号公式}是由真值和语句组成的一个有序对，即要么是：
  \[
  \sFmla{\True}{!A} \text{ 要么是 } \sFmla{\False}{!A}.
  \]
\end{defn}

直观上，我们可以把 $\sFmla{\True}{!A}$ 读作“$!A$ 可能为真”，把 $\sFmla{\False}{!A}$ 读作“$!A$ 可能为假”（在某个结构中）。

树中的每个带号公式要么是一个\emph{假设}（列在树的最顶端），
要么是上方的某个带号公式通过多条推理规则中的某一条得到的。
对于上方公式的每种可能的主联结词，都有两条规则：
一条用于标记为~$\True$ 的情况，另一条用于标记为~$\False$ 的情况。
有些规则允许树分叉，有些则只在枝上添加带号公式。
一条规则可以（而且常常必须）不是应用于紧邻其上的带号公式，
而是应用于该枝中从根节点到规则应用位置之间的任意带号公式。

当一个枝同时包含 $\sFmla{\True}{!A}$ 和 $\sFmla{\False}{!A}$ 时，该枝是\emph{封闭的}。
封闭的分析表是指每个枝都已封闭的分析表。
在直观解释下，任何一个枝都描述了一种联合的可能性，
但 $\sFmla{\True}{!A}$ 和 $\sFmla{\False}{!A}$ 是不可能联合出现的。
换句话说，如果一个枝是封闭的，它所描述的可能性就被排除了。
特别地，这意味着封闭的分析表排除了如下所有可能性：
同时使得每个形如 $\sFmla{\True}{!A}$ 的假设为真，且每个形如~$\sFmla{\False}{!A}$ 的假设为假。

\emph{针对 $!A$ 的封闭分析表}是以~$\sFmla{\False}{!A}$ 为根的封闭分析表。
如果这样的封闭分析表存在，$!A$ 为假的所有可能性就都已被排除；
也就是说，$!A$ 在每个结构中必然为真。

\end{document}